\section{Running the pipeline}
\label{sec:pipeline}

Eight notebooks run in the order in Table~\ref{tab:stages}, which is not the alphabetical order the
folder shows. Each stage reads the previous stage's files from disk and writes its own; nothing is
passed in memory, which is what makes the sequence safe to stop and restart.

\begin{table}[htbp]
\centering
\caption{The eight stages, in run order, with what to expect on time.}
\label{tab:stages}
\small
\begin{tabular}{|c|p{3.2cm}|p{4.6cm}|p{2.6cm}|p{2.4cm}|}
\hline
\textbf{\#} & \textbf{Notebook} & \textbf{Writes} & \textbf{From nothing} & \textbf{On this
folder} \\
\hline
1 & steam\_catalogue & the screened candidate list & 45 min & not needed \\
\hline
2 & review\_collection & a review file and checkpoint per game & 2 h 15 m & skipped \\
\hline
3 & auxiliary\_sources & per-game metadata & a few min & under a min \\
\hline
4 & review\_cleaning & cleaned English reviews & 30 min & under a min \\
\hline
5 & toxicity\_scoring & toxicity and sentiment scores & 1 to 2 h on GPU & under a min \\
\hline
6 & feature\_engineering & one row per game & 10 min & 10 min \\
\hline
7 & health\_score\_model & scores, bands, result tables & 2 to 3 min & 2 to 3 min \\
\hline
8 & figures & the eleven charts & 1 min & 1 min \\
\hline
\end{tabular}
\end{table}

\textbf{Stages 1 to 3 do not need running.} Their output ships in the folder and they account for
almost all the time. Starting at stage 4 reproduces every number in the research report. Run them
only to collect a fresh corpus, in which case the results will differ, because Steam has moved on
since 1 August 2026.

\textbf{To run a stage:} open the notebook, confirm the kernel reads \textbf{Python 3.11 (EA
Health)}, then choose \textbf{Restart Kernel and Run All Cells}. Do not run cells individually and
out of order, which is the one habit that produces plausible wrong numbers with no error. Every
notebook ends with an assertion cell checking its own output; if that cell prints its summary
without raising, the stage finished.

\subsection*{Stage 1: the candidate list}

Enumerates every multiplayer game on the store and screens each for an early-access history and
enough reviews to measure. Of 21,330 titles it keeps 1,887.

\subsection*{Stage 2: collecting the reviews}

The long one: 628,389 requests over about two and a quarter hours with eight workers. It is
rate-limited and backs off when the endpoint asks, so do not raise the worker count. It is also
resumable, because the checkpoint is written before the review file is renamed into place. The
shipped corpus was in fact collected across an interruption.

\shotpair{catalogue_screen}{The stage 1 screening result. Most of the catalogue is rejected for
too few reviews.}{collection_run}{Collection in progress, one game per line. A game already on
disk is skipped.}


\textbf{Trap.} Steam's own \texttt{app\_release\_date} is not the date a game left early access; it
differs by over a week for 74.8 per cent of launched titles. Stage 2 derives the real transition by
binary search over the review stream, and every later stage uses that derived date.

\subsection*{Stages 3 and 4: metadata and cleaning}

Stage 3 fetches app details, skipping anything already on disk. Stage 4 strips markup and links,
drops non-English reviews, hashes the author identifier, and keeps Steam's profanity mask as a
countable token rather than deleting it. From 3,504,533 collected reviews it keeps 1,865,234
English ones. Games already cleaned are skipped.

\shotpair{checkpoint_json}{One checkpoint. Besides the counts it records the derived launch date
and the binary-search probes the later stages read.}{cleaning_run}{Cleaning, one game per line.
Dropped rows are counted, never silently discarded.}

\subsection*{Stage 5: toxicity and sentiment scoring}

The stage that needs the graphics card. Detoxify scores toxicity, VADER scores sentiment, and
1,512,547 rows are written. The first run downloads the model weights once. Scoring is capped at
2,000 reviews per game, spread across eight time periods within each of the early-access and
post-launch phases so a quiet stretch keeps its representation; the cap binds on 318 of 1,793
games.

\noindent The first cell of this stage prints whether the card is visible. \texttt{cuda available :
False} still runs, on the processor, in hours rather than minutes. The resulting score distribution
skews hard to zero, as a healthy result should: about three per cent of reviews score above 0.5.

\subsection*{Stages 6 to 8: features, models and figures}

Stage 6 reduces the scored reviews to one row per game, 1,793 rows and 35 columns, and computes the
retention measure as the gap between a reviewer's lifetime playtime and their playtime when
writing. Stage 7 narrows to the 545 games with enough reviews either side of launch, splits them by
launch date into 381 training and 164 test, fits the models and writes the nineteen result tables.
Stage 8 draws the eleven charts.

Everything fitted, including the scaling and the band cutoffs, is fitted on the training partition
only, and the test partition is scored once at the end.

\shotpair{feature_table}{The finished feature table, one row per
game.}{model_comparison}{The results table as the notebook prints it. The baseline row is the
number the project has to beat; the gain column is the claim the research report makes.}

\shotpair{assertion_cell}{The closing assertion cell, re-checking that the split really was
temporal before declaring the stage complete.}{figure_inline}{A chart rendered in the notebook.
Each is also written to \texttt{figures} at 300 dots per inch.}
